%%% MA_08_GE_ISLM.tex -- reorganised; frames moved from the old MA_xx files (see comments)

\section{Fiscal Policy, General Equilibrium and the IS-LM Model}

%%\section{The Goods Market: Aggregate Demand and the Multiplier}

%%% [moved from MA_5_General_Equlibrium.tex, frame 1]
\begin{frame}{Say's Law}
    \mytext{
    Every demand equals supply, as we produce to have something to exchange for what we desire.

    Every supply equals demand, as everything produced will at some point be exchanged (or self-consumed).

    Production is the source of demand.
    }
\end{frame}

%%% [moved from MA_1_Fundamentals.tex, frame 23]
 

\begin{frame}{Keynesian Cross}
    \showmytikz{../shared/tikz/tikz-keynesiancross}
    \mycomment{
    Aggregate demand $AD$ is consumption, investment, and government expense;
    \begin{equation}
        AD = C+I+G
    \end{equation}
    with
    \begin{equation}
        C=a+b\left(Y-t\right)
    \end{equation}
    and $I,G,t$ constant; so that
    \begin{equation}
        AD = e + b\cdot Y
    \end{equation}
    with $G$ part of the autonomous part $e$.
    }
\end{frame}

%%% [moved from MA_1_Fundamentals.tex, frame 24]
\begin{frame}{Keynesian Cross}
    \showmytikz[2]{../shared/tikz/tikz-keynesiancross}
    \mycomment{
    An increase of government spending increases aggregate demand, shifting the $AD$ curve up. The result is a higher output $Y$.

    The same applies for autonomous increase of $I$.
    }
\end{frame}

%%% [moved from MA_1_Fundamentals.tex, frame 25]
\begin{frame}{Fiscal Multiplier}
    \mycomment{
    The idea with the multiplier is that any extra inserted expenditure (by fiscal policy, but anything else for that matter) is someone else's income, who will then re-spend a part of that, saving another part. 

    The size of the multiplier will depend on the propensity to spend (the inverse of the propensity to save), and on taxes paid each time.
    }
    \showmytikz{../shared/tikz/tikz-multiplier}
\end{frame}

%%% [moved from MA_1_Fundamentals.tex, frame 26]
\begin{frame}{Fiscal Multiplier}
    \mycomment{
    For now, let us abstract from taxes and net exports.
    }
    \mytext{
    Consumption $C$ is a (constant) part of income $Y$:
    \begin{equation}
        C=C_0+c\cdot Y
    \end{equation}
    Thus
    \begin{equation}
        Y=C_0+c\cdot Y + I + G
    \end{equation}
    For (in the short run) constant $I,C_0$ we get
    \begin{equation}
        \Delta Y = \frac{\Delta G}{1-c}=\frac{1}{s}
    \end{equation}
    with $s$ the savings quota. --
    Thus the multplier is
    \begin{equation}
        M=\frac{1}{1-c}=\frac{1}{s}.
    \end{equation}
    }
    \mycomment{
    % There are a number of (simplifying) assumptions here. One is that net exports $NX=Ex-Im$ be zero. Sin
    }
\end{frame}

%%% [moved from MA_1_Fundamentals.tex, frame 27]
\begin{frame}{Fiscal Multiplier: Critique}
    \mytext{
    \begin{itemize}
        \item crowding out private demand
        \item government debt means higher taxes:
        \item[$\rightarrow$] subjects reduce spending
        \item debt financing raises interest rates
        \item monetary policy counteracts
        \item lock-in problem
        \item might depend on government debt level
    \end{itemize}
    }
\end{frame}


%%% TODO (new content): automatic stabilisers; tax multiplier / balanced-budget multiplier; open-economy multiplier (import leakage)

%%%%%

%%% [moved from MA_5_General_Equlibrium.tex, frame 0]
\begin{frame}{From Partial to General Equilibrium}
    \mytext{
    Microeconomics demonstrated the equilibrium of supply and demand for individual markets: partial equilibria.

    We have also seen that goods can be complementary or substitutes. The price of one good (equilibrium on that market's good) influences the price of another (cross-price elasticities).

    So the questions are: how does one market influence another? What can we say about an equilibrium on multiple markets, a general equilibrium?
    }
\end{frame}

%%% [moved from MA_5_General_Equlibrium.tex, frame 2]
\begin{frame}{Walras' Law}
    \mytext{
        If one market is in disequilibrium, another market must also be in disequilibrium.

        The sum of all derivations from equilibrium is zero.

        It follows that if $n-1$ markets are in equilibrium, market $n$ is also in equilibrium.
    }
\end{frame}

%%% [moved from MA_5_General_Equlibrium.tex, frame 3]
\begin{frame}{Patinkin's Model of Money, Interest, and Prices}
    \showmytikz[-1,1]{../shared/tikz/tikz-threemarkets}
    \mycomment{
    Think of two (partly substitutable) goods $X,Y$ and a denominator $Z$, like fish, coconut and shells. 

    As the price of one good increases, the price of the other increases as well (positive cross-price effect, substitutes). 

    There is a price ratio for $X,Y$  where supply and demand for $X$ are in equilibrium.

    \bigskip 
    
        All this follows Krugman: \href{https://web.mit.edu/krugman/www/islm.html}{https://web.mit.edu/krugman/www/islm.html}
    
    }
\end{frame}

%%% [moved from MA_5_General_Equlibrium.tex, frame 4]
\begin{frame}{Patinkin's Model of Money, Interest, and Prices}
    \showmytikz[-2,1]{../shared/tikz/tikz-threemarkets}
    \mycomment{
    The same applies for $Y$. 
    
    Both slope upwards since they are substitutes; their cross-price elasticity is positive.
    
    Now in the case of only two goods, both markets must be in equilibrium simultaneously. Thus, there will be a certain number  $X^*,Y^*$ sold and bought for a specific price.
    }
\end{frame}

%%% [moved from MA_5_General_Equlibrium.tex, frame 5]
\begin{frame}{Patinkin's Model of Money, Interest, and Prices}
    \showmytikz[1]{../shared/tikz/tikz-threemarkets}
    \mycomment{

    
    Walras' law states: when $n-1$ markets are in balance, the remaining market must also be in balance.

    Thus, we can also have a third good $Z$, in our example the shells. 
    
    % Think of two (partly compementary) goods $X,Y$ and a denominator $Z$, like fish, coconut and shells.
    }
\end{frame}

%%% [moved from MA_5_General_Equlibrium.tex, frame 6]
\begin{frame}{Patinkin's Model of Money, Interest, and Prices}
    \showmytikz[2]{../shared/tikz/tikz-threemarkets}
    \mycomment{
    We can think of these three goods as goods `in general' ($X$), bonds `in general' ($Y$) and money ($Z$). (Or consumption vs. savings.)

    So money can be either held, or spent on goods (consumed), or invested in bonds.

    The price of bonds is typically expressed by the interest rate; so the price axis `price of bonds' is now inverse. Let us flip it so that the axis shows the interest rate.
    }
\end{frame}

%%% [moved from MA_5_General_Equlibrium.tex, frame 7]
\begin{frame}{Patinkin's Model of Money, Interest, and Prices}
    \showmytikz[3]{../shared/tikz/tikz-threemarkets}
    \mycomment{
    Now we have a general equilibrium model of three markets: goods, money, and bonds.
    }
\end{frame}

%%% [moved from MA_5_General_Equlibrium.tex, frame 8]
\begin{frame}{Patinkin's Model of Money, Interest, and Prices}
    \showmytikz[3,4]{../shared/tikz/tikz-threemarkets}
    \mycomment{
        If the market for goods remains stable, but there is more supply of money, interest rates will go down and the price of goods will increase.

        This is equal to a monetary stimulus: the government (central bank) buys bonds and `sells' newly printed money. Cheper money reduces the interest rate at higher output.
    }
\end{frame}

%%% [moved from MA_5_General_Equlibrium.tex, frame 9]
\begin{frame}{Patinkin's Model of Money, Interest, and Prices}
    \showmytikz[3,5]{../shared/tikz/tikz-threemarkets}
    \mycomment{
        If money supply remains stable, an increase of interest rates will result in a higher price for goods.

        Or say it this way: an investment boom (bonds curve shift) leads to higher output (goods curve shift) at higher interest rates ($i_2>i^*$) and higher inflation ($P_2>P^*$).

        This will also happen when the government sells bonds and buys goods (fiscal stimulus).
    }
\end{frame}

%%% [moved from MA_5_General_Equlibrium.tex, frame 10]
% \begin{frame}{Patinkin's Model of Money, Interest, and Prices}
%     \showmytikz[3,6]{../shared/tikz/tikz-threemarkets}
%     \mycomment{
%         A stable supply of bonds means more money will push the price of goods down at lower interest rates. 
        
        
%     }
% \end{frame}

%%% [moved from MA_6_ISLM.tex, frame 0]
\begin{frame}{IS-LM Model: IS Curve}
    \showmytikz[1,2,2.5]{../shared/tikz/tikz-islm}
    \mycomment{
   $IS$ is the curve of all interest rate - output - combinations that result in a balance of the real markets, that is, of investments~$I$ and savings~$S$.

    At $P_1$, a relatively high interest rate makes credit expensive. At this rate, much will be saved, hence less will be spent on consumption. There will be relatively low output.

    At $P_2$, lower interest rate make consumption more interesting than saving. Thus, output $Y$ will be higher.

    % Empirically, the $IS$ curve is relatively steep, while the $LM$ curve is relatively flat. 
    }
\end{frame}

%%% [moved from MA_6_ISLM.tex, frame 1]
\begin{frame}{IS-LM Model: LM Curve}
    \showmytikz[3,4,4.5]{../shared/tikz/tikz-islm}
    \mycomment{
        $LM$ is the curve of all combinations of interest rate and output that result in a balance on the money market, namely of liquidity preference $L$ and money supply $M$.

        At $P_3$, a relatively low output will require little money, and low money demand will result in low interest rates. 

        At $P_4$, a higher output will require more money, and as more profits are realised, interest rates will rise.

        The exact shape of the $LM$ curve is less clear than that of the $IS$ curve. In reality, it might be very flat. Some say it is flat on the left (liquidity trap) and steep on the right.        
        
        % Future expectations and policy play a larger role on the financial markets than in the real markets. 
    }
\end{frame}

%%% [moved from MA_6_ISLM.tex, frame 2]
\begin{frame}{IS-LM Model: Equilibrium}
    \showmytikz[1,3,5]{../shared/tikz/tikz-islm}
    \mycomment{
        There is an equilibrium $A$ where $IS$ and $LM$ meet, that is the synchronous balance of real and financial markets. At this combination of interest rate $i$ and output $Y$ investments will equal savings (the point is on $IS$) and money supply will equal money demand (the point is on $LM)$.
    }
\end{frame}

%%% [moved from MA_6_ISLM.tex, frame 3]
\begin{frame}{IS-LM Model: Disequilibria}
    \showmytikz[1,2,3,4,5]{../shared/tikz/tikz-islm}
    \mycomment{
        $P_1$: high interest rates result in little spending and low output. At this interest rate there is an excess of money supply (the point is left of $LM$). Money suppliers will lower the interest rate to find clients, which will result in a higher output, so movement towards $A$.

        $P_2$: at a low interest rate and high output, there is excess demand for money for transactions. This will drive up interest rates. 

        $P_3$: low interest rate and low output means there is much cheap money for more consumption available, creating excess demand for goods, driving up output. 

        $P_4$: high interest rates clear the money market, but the output is too high at this point for a balanced real market; subjects would rather save than consume this much. Thus, production and interest rates will have to go down.
        
    }
\end{frame}

%%% [moved from MA_6_ISLM.tex, frame 4]
\begin{frame}{IS-LM Model: Fiscal Policy}
    \showmytikz[1,3,5,6]{../shared/tikz/tikz-islm}
    \mycomment{
        Government expenditure can shift the $IS$ curve to the right, resulting in higher interest rate $i_2>i_1$ and larger output $Y_2>Y_1$. 
        
        The effect is only partial; the rise in output is smaller than the additional government expenditure. The risen interest rate has resulted a partial crowding-out. 

        Tariffs have a similar effect, as they raise net exports ($NX=EX-IM$). 

        Any fiscal policy has also distributive effects not shown here.
    }
\end{frame}

%%% [moved from MA_6_ISLM.tex, frame 5]
\begin{frame}{IS-LM Model: Monetary Policy}
    \showmytikz[1,3,5,7]{../shared/tikz/tikz-islm}
    \mycomment{
        Monetary policy (by the central bank) can shift the $LM$ curve to the right. This can happen either through lowering of the interest rate, or by open market policy. 

        The new equilibrium shows a lower interest rate and higher output $Y_3>Y_1$. The effect is also only partial: the rise in output is smaller than the monetary stimulus, since the lower interest rate has reduced savings and investment.

        There is an obvious contradiction (or: competition?) between fiscal and monetary policy, which is the reason why central banks \textit{should be} independent.
    }
\end{frame}

%%% [moved from MA_6_ISLM.tex, frame 6]
\begin{frame}{IS-LM Model: Fiscal vs. Monetary Policy}
    \showmytikz[1,3,5,6,7]{../shared/tikz/tikz-islm}
    \mycomment{
        Fiscal and monetary policy can complement or offset each other. Expansion by policy will meet its natural boundary at unsustainable government debt levels or at runaway inflation discrediting the currency. 
        
        It is thus widely believed that any such policy should be used temporarily only, like to overcome an acute crisis, external shocks, slumps, or maybe to catch up on development. 
        
        It is a bit like burning furniture to heat the room, or selling the silverware: you can do it in a crisis, but it is not a long-term solution.
    }
\end{frame}

%%% [moved from MA_6_ISLM.tex, frame 7]
   
\begin{frame}{IS-LM Model, Curved: Fiscal Policy}
    \showmytikz[1,3,5,6]{../shared/tikz/tikz-islm-smooth}
    \mycomment{
        The actual shape of IS and LM are debated. They `seem' to be steep at the extremes, but flat for low levels of $i$. This can, of course, have dramatic consequences whether fiscal or monetary policy works at all. 

        A steep $IS$-curve indicates that fiscal policy has only little impact on interest rates. Then there is relatively little crowding-out.
    }
\end{frame}

%%% [moved from MA_6_ISLM.tex, frame 8]
    \begin{frame}{IS-LM Model, Curved: Monetary Policy}
    \showmytikz[1,3,5,7]{../shared/tikz/tikz-islm-smooth}
    \mycomment{
        The shape of the $LM$-curve is much disputed. The consensus sees it flat for low $Y,i$ and steep for higher $Y,i$. 

        In the lower region, the flat $LM$-curve implies effectivelessness of monetary policy; larger money supply does not lead to more investment, since economic subjects have a strong preference for liquidity over investing. This situation is called `liquidity trap'.
    }
\end{frame}

